% 第 2 章 原子核结构理论模型与方法

\chapter{原子核结构理论模型与集体运动理论框架}
\label{ch:theoretical_models}

\section{单粒子运动与变形 Nilsson 平均场}
\label{sec:nilsson_model}

要透彻理解原子核为什么会发生转动，首先必须理解单个核子在变形平均场中是如何运动的。在球对称势阱中，单核子轨道完全由主量子数 $N$、轨道角动量 $l$、总角动量 $j$ 以及磁量子数 $m$ 标记，具有 $(2j+1)$ 重磁简并。

当四极相互作用将原子核拉伸为一个椭球体时，球对称性荡然无存。S. G. Nilsson（1955）在各向异性谐振子势的基础上，构建了经典的变形单粒子壳模型哈密顿量：
\begin{equation}
    h_{\text{Nilsson}} = -\frac{\hbar^2}{2m} \nabla^2 + \frac{1}{2} m \left( \omega_x^2 x^2 + \omega_y^2 y^2 + \omega_z^2 z^2 \right) - \kappa \hbar \mathring{\omega}_0 \left[ 2 \vec{l} \cdot \vec{s} + \mu (\vec{l}^2 - \langle \vec{l}^2 \rangle_N) \right]
    \label{eq:nilsson_hamiltonian}
\end{equation}
其中，体积守恒条件要求 $\omega_x \omega_y \omega_z = \mathring{\omega}_0^3 = \text{const}$。

在柱坐标与拉伸坐标系 $(\xi, \eta, \zeta)$ 下，各向异性频率可参数化为四极形变参数 $\epsilon$（或 $\beta_2$）：
\begin{equation}
    \omega_z = \mathring{\omega}_0 \left( 1 - \frac{2}{3} \epsilon \right), \quad \omega_\perp = \omega_x = \omega_y = \mathring{\omega}_0 \left( 1 + \frac{1}{3} \epsilon \right)
    \label{eq:nilsson_frequencies}
\end{equation}
式中自旋-轨道耦合项 $2\vec{l}\cdot\vec{s} = \vec{j}^2 - \vec{l}^2 - \vec{s}^2$ 负责产生球形幻数；而扁平修正项 $\mu (\vec{l}^2 - \langle \vec{l}^2 \rangle_N)$（其中 $\langle \vec{l}^2 \rangle_N = \frac{1}{2}N(N+3)$）用于压低高角动量轨道的能量，以弥补谐振子势阱与真实平底 Woods-Saxon 势之间的形状偏差。

在非轴对称三轴椭球下，三个主轴频率由四极形变参数 $\beta_2$ 与三轴不对称度 $\gamma$ 严格解耦：
\begin{equation}
    \omega_k = \mathring{\omega}_0 \exp\left[ -\sqrt{\frac{5}{4\pi}} \beta_2 \cos\left(\gamma - \frac{2\pi}{3} k\right) \right], \quad k = 1, 2, 3
    \label{eq:triaxial_frequencies}
\end{equation}
频率的体积守恒条件自然满足：$\prod_{k=1}^3 \omega_k = \mathring{\omega}_0^3$。三轴形变项使得单核子势能面中引入了非对角 $\Delta \Omega = \pm 2$ 的四极耦合，使得单粒子轨道不再保持任何单轴投影守恒，促成空间中不同磁子态的强相干杂化。

在轴对称椭球极限下（$\omega_x = \omega_y \ne \omega_z$），总角动量 $j$ 不再守恒，只有总角动量在核对称轴（$z$ 轴）上的投影 $\Omega$ 保持为良好量子数。每一个球形 $j$ 轨道因此分裂为 $(2j+1)/2$ 对互为简并的能量能级（$\pm \Omega$ 对，Kramers 简并）：
\begin{equation}
    \Omega = 1/2, 3/2, \dots, j
\end{equation}

在大形变极限下，单核子状态通常由一组直观的\term{渐近量子数} $[N n_z \Lambda] \Omega$ 标记：
\begin{itemize}
    \item $N$：主振动态量子数；
    \item $n_z$：核子波函数沿对称轴（$z$ 轴）振动的节点数；
    \item $\Lambda$：轨道角动量在对称轴上的投影（$\Omega = \Lambda \pm 1/2$）。
\end{itemize}

初学者应当建立极其清晰的几何直觉：$n_z$ 代表核子波函数沿拉长方向的空间振荡密度。$n_z$ 越大的轨道，核子在空间上越集中分布在椭球两极的拉长尖端；当原子核发生长椭形变时，核子与周围核物质重叠最深，引力势阱最深，因此能级随形变 $\beta$ 增加而\textbf{急剧下沉}（图~\ref{fig:nilsson_splitting}）。相反，$n_z = 0$ 的轨道核子主要分布在赤道平面上，受排斥作用能量随形变升高。

\begin{figure}[htbp]
    \centering
    \begin{tikzpicture}[scale=0.95]
        \draw[thick, ->] (0,0) -- (6.5,0) node[right] {形变参数 $\beta_2$};
        \draw[thick, ->] (0,-3) -- (0,4) node[above] {单粒子能量 $\epsilon$};

        % 球形能级
        \node[left] at (0,1.5) {$1h_{11/2}$};

        % 分裂曲线
        \draw[line width=1.2pt, red] (0,1.5) .. controls (2,0.8) and (4,-1.5) .. (6,-2.8) node[right] {$[550]1/2^-$ (强下沉)};
        \draw[line width=1.0pt, red!80!black] (0,1.5) .. controls (2,1.0) and (4,-0.5) .. (6,-1.2) node[right] {$[541]3/2^-$};
        \draw[line width=1.0pt, orange] (0,1.5) .. controls (2,1.3) and (4,0.5) .. (6,0.2) node[right] {$[532]5/2^-$};
        \draw[line width=1.0pt, blue!80!black] (0,1.5) .. controls (2,1.6) and (4,1.8) .. (6,1.8) node[right] {$[523]7/2^-$};
        \draw[line width=1.0pt, blue] (0,1.5) .. controls (2,1.9) and (4,2.8) .. (6,3.0) node[right] {$[514]9/2^-$};
        \draw[line width=1.2pt, purple] (0,1.5) .. controls (2,2.2) and (4,3.5) .. (6,3.9) node[right] {$[505]11/2^-$ (强上升)};
    \end{tikzpicture}
    \caption{高 $j$ 闯入轨道 $1h_{11/2}$ 随四极形变 $\beta_2$ 产生能级分裂示意图}
    \label{fig:nilsson_splitting}
\end{figure}

特别值得注意的是所谓\term{高 $j$ 闯入轨道}（Intruder Orbitals，如 A~80 区的 $1g_{9/2}$、A~130 区的 $1h_{11/2}$ 以及稀土区的 $1i_{13/2}$）。这些轨道因自旋-轨道耦合向下跨越整整一个主壳层，其宇称与周围正常轨道完全相反，无法被同壳层轨道混合阻尼。处于壳层底部的最低 $\Omega$ 轨道（如 $[550]1/2^-$）随形变极深下潜，成为驱动整个原子核发生显著形变的发动机。

\section{配对相互作用与 BCS 准粒子概念}
\label{sec:bcs_pairing}

如果原子核仅受变形单粒子势阱约束，核子将按 Pauli 原理简单从低到高填满能级。然而，核子之间存在着强大的短程吸引残余相互作用——\term{配对相互作用}（Pairing Interaction）。它使两个处于互为时间反演轨道 $|k\rangle$ 与 $|\bar{k}\rangle$ 上的全同核子倾向于相干凝聚为自旋为零的超流玻色子对（Cooper 对）。

在简单单极配对力模型下，相互作用哈密顿量表达为：
\begin{equation}
    H_{\text{pair}} = -G \sum_{k, k'>0} c_k^\dagger c_{\bar{k}}^\dagger c_{\bar{k}'} c_{k'} = -G \hat{P}^\dagger \hat{P}
    \label{eq:monopole_pairing}
\end{equation}
其中 $G$ 为配对相互作用强度参数（中重核通常 $G \approx 20/A\text{ MeV}$），对产生算符 $\hat{P}^\dagger = \sum_{k>0} c_k^\dagger c_{\bar{k}}^\dagger$ 将一对处于共轭时间反演态上的粒子作为一个整体产生。

在 Bardeen-Cooper-Schrieffer（BCS）理论框架下，核子算符通过 Bogoliubov 变换改写为\term{准粒子算符}（Quasiparticle Operators）：
\begin{equation}
    \alpha^\dagger_k = u_k c^\dagger_k - v_k c^\dagger_{\bar{k}}, \quad \alpha^\dagger_{\bar{k}} = u_k c^\dagger_{\bar{k}} + v_k c^\dagger_k
    \label{eq:bogoliubov_trans}
\end{equation}
其中 $v_k^2$ 代表单粒子能级 $k$ 的占据几率，$u_k^2 = 1 - v_k^2$ 代表未占据几率：
\begin{equation}
    v_k^2 = \frac{1}{2} \left[ 1 - \frac{\epsilon_k - \lambda}{\sqrt{(\epsilon_k - \lambda)^2 + \Delta^2}} \right]
\end{equation}
这里 $\lambda$ 为体系费米能级（Fermi Surface），$\Delta$ 为微观配对能隙（Pairing Gap Parameter）。

配对能隙 $\Delta$ 与平均核子数约束方程通过变分法自洽确定（BCS 能隙方程）：
\begin{equation}
    \Delta = G \sum_{k>0} u_k v_k = \frac{G}{2} \sum_{k>0} \frac{\Delta}{\sqrt{(\epsilon_k - \lambda)^2 + \Delta^2}}
    \label{eq:gap_equation}
\end{equation}
\begin{equation}
    N = 2 \sum_{k>0} v_k^2 = \sum_{k>0} \left[ 1 - \frac{\epsilon_k - \lambda}{\sqrt{(\epsilon_k - \lambda)^2 + \Delta^2}} \right]
    \label{eq:number_equation}
\end{equation}

单准粒子激发能量为：
\begin{equation}
    E_k = \sqrt{(\epsilon_k - \lambda)^2 + \Delta^2}
\end{equation}

对于初学者，BCS 理论带来了两个具有颠覆性的物理启示：
\begin{enumerate}
    \item \textbf{激发能谱的能隙存在}：在偶偶核中，打破一对 Cooper 对至少需要消耗 $2\Delta \approx 2 \text{ MeV}$ 的能量。这完美解释了为什么所有偶偶核的第一个内禀准粒子激发态都在 1.5 MeV 以上，而其下方只有纯粹的基态集体转动谱；
    \item \textbf{奇核的阻塞效应（Blocking Effect）}：在奇 $A$ 核中，未配对的奇价核子牢牢占据了特定轨道 $k_0$，由于 Pauli 原理禁止其他核子进入，该轨道完全退出了对凝聚过程。这直接导致奇核中的有效配对能隙明显缩小（通常 $\Delta_{\text{odd}} \approx 0.8 \Delta_{\text{even}}$），转动惯量显著增大。
\end{enumerate}

\section{推转壳模型（CSM）与高自旋能带动力学}
\label{sec:cranked_shell_model}

\subsection{转动坐标系与 Routhian 能量}

当原子核以角频率 $\omega$ 绕某一固定轴（如 $x$ 轴）高速旋转时，实验室系中的时间依赖哈密顿量极其复杂。R. Bengtsson 与 S. Frauendorf（1979）引入了著名的\term{推转壳模型}（Cranked Shell Model, CSM）。在随动转动坐标系中，体系哈密顿量变换为非时间依赖的\term{推转 Routhian}（Routhian）：
\begin{equation}
    H' = H - \omega J_x = \sum_k (h_{\text{Nilsson}} - \lambda \hat{N}) - \Delta (P^\dagger + P) - \hbar\omega J_x
    \label{eq:csm_routhian}
\end{equation}

推转项 $-\omega J_x$ 在经典力学上对应着转动系中的\textbf{科里奥利力（Coriolis Force）}。Coriolis 力正比于核子的角动量及其线速度：它试图将每一个独立运动核子的内禀轨道角动量强行拉向转动轴（$x$ 轴）。这彻底破坏了时间反演不变性，原先在静态平均场中简并的对核子发生能量分道。

在两准粒子表象中，对哈密顿量进行 Bogoliubov-de Gennes 对角化：
\begin{equation}
    \begin{pmatrix} h - \lambda - \omega j_x & \Delta \\ \Delta & -(h - \lambda) - \omega j_x \end{pmatrix} \begin{pmatrix} U_i \\ V_i \end{pmatrix} = e'_i \begin{pmatrix} U_i \\ V_i \end{pmatrix}
    \label{eq:hfb_cranking}
\end{equation}
特征值 $e'_i(\omega)$ 即为准粒子在转动参考系中的准粒子 Routhian 能级。根据热力学广义位移关系，准粒子对转动轴贡献的顺磁对齐角动量（Alignment）直接等于准粒子能量对转速的负斜率：
\begin{equation}
    i_x = -\frac{d e'_i}{d\omega} = \langle i | j_x | i \rangle
    \label{eq:alignment_formula}
\end{equation}

\subsection{准粒子顺磁排列与能带回弯微观机制}

随着转动频率 $\hbar\omega$ 的不断提高，处于高 $j$ 闯入轨道（如 $1h_{11/2}$）上的最靠近费米面的核子对，所感受到的 Coriolis 去耦力越来越强。在某一临界频率 $\hbar\omega_c$ 处，打破这对核子的能耗被 Coriolis 对齐能完全抵消，两个准粒子的角动量同时顺磁对齐到转动轴上，贡献了巨大的角动量增量（可达 $i_x \approx 8 \sim 10\hbar$）。

对齐带来的角动量突增，在实验绘制的转动惯量与推转频率曲线（$\mathcal{J}^{(1)}$ vs $\hbar\omega$）上表现为剧烈的反折，因形似弯腰回折而被称为\term{回弯}（Backbending），如图~\ref{fig:backbending_interaction}所示。

\begin{figure}[htbp]
    \centering
    \begin{tikzpicture}[scale=0.9]
        % 左图：尖锐回弯
        \begin{scope}[shift={(0,0)}]
            \draw[thick, ->] (0,0) -- (4.5,0) node[right] {$\hbar\omega$ (MeV)};
            \draw[thick, ->] (0,0) -- (0,3.5) node[above] {$\mathcal{J}^{(1)}$};
            \draw[line width=1.3pt, blue] (0.5,0.5) -- (2.0,1.2) .. controls (2.6,1.8) and (2.7,2.2) .. (2.3,2.5) .. controls (2.0,2.7) and (2.2,3.0) .. (3.8,3.2);
            \node at (2.2,-0.6) {(a) 尖锐回弯 ($|V_{12}| < 35$ keV)};
        \end{scope}

        % 右图：平滑上弯
        \begin{scope}[shift={(6.5,0)}]
            \draw[thick, ->] (0,0) -- (4.5,0) node[right] {$\hbar\omega$ (MeV)};
            \draw[thick, ->] (0,0) -- (0,3.5) node[above] {$\mathcal{J}^{(1)}$};
            \draw[line width=1.3pt, red!80!black] (0.5,0.5) .. controls (2.0,1.0) and (2.2,1.8) .. (3.8,3.0);
            \node at (2.2,-0.6) {(b) 平滑上弯 ($|V_{12}| \ge 65$ keV)};
        \end{scope}
    \end{tikzpicture}
    \caption{两能带交叉相互作用强度 $|V_{12}|$ 控制回弯形态的物理机制}
    \label{fig:backbending_interaction}
\end{figure}

在转动参考系两能带混合模型中，基态带（$g$-带）与对齐带（$s$-带，Stockholm 带）在回避交叉点处的相互作用哈密顿量为：
\begin{equation}
    H'(\omega) = \begin{pmatrix} e'_g(\omega) & V_{12} \\ V_{12} & e'_s(\omega) \end{pmatrix}
\end{equation}
交叉点两本征带的最小能量间距严格等于两倍相互作用矩阵元：$\Delta e'_{\text{min}} = 2 |V_{12}|$。
\begin{itemize}
    \item 当 $|V_{12}|$ 很小（$< 35\text{ keV}$）时，两带在交叉点几乎不混合，转动能谱产生尖锐回弯（$S$ 形反折）；
    \item 当 $|V_{12}|$ 很大（$\ge 65\text{ keV}$）时，两带发生强烈混合，能级平滑过渡，仅在能谱上表现为动力学转动惯量的平缓上升，即所谓的\term{平滑上弯}（Smooth Upbending）。
\end{itemize}

\subsection{旋称量子数（Signature）与旋称反转}

在绕 $x$ 轴推转后，空间宇称 $\mathcal{P}$ 仍然守恒，但角动量在其他轴的投影已不守恒。体系此时对绕转动轴旋转 $180^\circ$ 的空间操作 $\mathcal{R}_x(\pi) = e^{-i \pi J_x}$ 保持对易守恒。对应的离散守恒对称性量子数称为\term{旋称}（Signature）：
\begin{equation}
    r = e^{-i \pi \alpha}, \quad \alpha = \begin{cases} 0, 1 & (\text{偶 } A \text{ 核}) \\ \pm 1/2 & (\text{奇 } A \text{ 核}) \end{cases}
\end{equation}

正常情况下，受 Coriolis 力推动，具有有利旋称（Favored Signature，如奇核中的 $\alpha = 1/2$）的转动分支在能量上永远低于不利旋称（Unfavored Signature，$\alpha = -1/2$）。然而，在 A~130 与 A~160 区域奇奇核与奇 A 核低自旋端，实验广泛观测到了不利分支反而处于低能量的异常现象，即\term{旋称反转}（Signature Inversion）。微观推导证明，这是原子核发生三轴形变（$\gamma < 0$）导致非对角张量耦合逆转的确定性微观证据。

\section{微观三轴转子理论与模型无关四极闭合}
\label{sec:triaxial_theory_and_closure}

\subsection{严格三轴粒子-转子模型（PRM）}

推转壳模型（CSM）本质上是一个半经典的平均场近似，角动量守恒性受到破坏。为了严格计算有限自旋下的离散态跃迁几率，微观核结构中建立了三轴粒子-转子模型（Particle-Rotor Model, PRM）：
\begin{equation}
    H_{\text{PRM}} = \sum_{k=1}^3 \frac{(I_k - j_k)^2}{2\mathcal{J}_k} + h_{\text{sp}}(\beta, \gamma) + H_{\text{pair}}
    \label{eq:prm_hamiltonian}
\end{equation}

PRM 将全系总角动量分解为内禀准粒子角动量 $\vec{j}$ 与变形核芯转动角动量 $\vec{R} = \vec{I} - \vec{j}$。通过构建包含完整三维 Wigner $D$-函数与单核子态的严格反对称化基矢基底 $|I M K \Omega\rangle$，进行完全微观对角化求解。该模型不仅能够严格保持角动量守恒，而且自然蕴含了摆动声子激发与手征几何退简并的所有量子穿透效应，是当前解释三轴集体运动的最权威基准工具。

哈密顿量展开后包含三部分：
\begin{equation}
    H_{\text{PRM}} = H_{\text{rot}}^{\text{core}} + H_{\text{Coriolis}} + H_{\text{recoil}} + h_{\text{sp}}
    \label{eq:prm_expanded}
\end{equation}
其中转子项为：
\begin{equation}
    H_{\text{rot}}^{\text{core}} = \sum_{k=1}^3 \frac{I_k^2}{2\mathcal{J}_k}
\end{equation}
科里奥利相互作用项为：
\begin{equation}
    H_{\text{Coriolis}} = -\sum_{k=1}^3 \frac{I_k j_k}{\mathcal{J}_k} = -\frac{1}{2} \left( \frac{1}{\mathcal{J}_1} + \frac{1}{\mathcal{J}_2} \right) (I_+ j_- + I_- j_+) - \frac{1}{4} \left( \frac{1}{\mathcal{J}_1} - \frac{1}{\mathcal{J}_2} \right) (I_+ j_+ + I_- j_-) - \frac{I_3 j_3}{\mathcal{J}_3}
    \label{eq:prm_coriolis}
\end{equation}
非轴对称因子 $\left( \frac{1}{\mathcal{J}_1} - \frac{1}{\mathcal{J}_2} \right)$ 明确诱发了 $\Delta K = \pm 2$ 与 $\Delta \Omega = \pm 2$ 的同时跳迁。这种强烈的张量耦合在内禀坐标系中将单准粒子角动量牢固拉入空间三维倾斜取向，成为从微观角度计算核手征性与摆动进动的严格定量工具。

基矢展开显式构造为：
\begin{equation}
    |I M \alpha\rangle = \sum_{K \Omega} C_{K \Omega}^\alpha \sqrt{\frac{2I+1}{16\pi^2 (1 + \delta_{K 0}\delta_{\Omega 0})}} \left[ D_{M K}^I(\alpha, \beta, \gamma) |\Omega\rangle + (-1)^{I-j} D_{M,-K}^I(\alpha, \beta, \gamma) |-\Omega\rangle \right]
    \label{eq:prm_basis}
\end{equation}
其中矩阵元中包含纯转子项、单粒子内禀项以及核心的非对角 Coriolis 交叉耦合项 $-\sum_k \frac{I_k j_k}{\mathcal{J}_k}$。通过全空间精确微观对角化，PRM 能够极其自然地计算出实验室系下的绝对电磁跃迁约化几率 $B(M1)$ 与 $B(E2)$，杜绝了半经典推转中角动量投影截断所引起的过渡伪影。

\subsection{三维倾斜轴推转模型（3D-TAC）与手征势能面}

经典推转模型（CSM）强行将转动轴限制在单核几何主轴（通常为 $x$ 轴）上，称为一维主轴推转（1D Principal Axis Cranking, PAC）。当原子核处于三轴非共面手征相态时，总角动量矢径 $\vec{J}$ 必然脱离主轴。S. Frauendorf 与孟杰（1997）发展了\term{三维倾斜轴推转模型}（3D Tilted Axis Cranking, 3D-TAC）。

在 3D-TAC 中，转动角速度矢量在球坐标下参数化为极角 $\theta$ 与方位角 $\phi$：
\begin{equation}
    \vec{\omega} = \omega (\sin\theta \cos\phi\, \vec{e}_x + \sin\theta \sin\phi\, \vec{e}_y + \cos\theta\, \vec{e}_z)
    \label{eq:3d_tac_omega}
\end{equation}
对应的推转 Routhian 为：
\begin{equation}
    h'(\theta, \phi) = h_{\text{Nilsson}}(\beta, \gamma) - \Delta (P^\dagger + P) - \lambda \hat{N} - \vec{\omega} \cdot \vec{j}
    \label{eq:3d_tac_routhian}
\end{equation}
自洽总能量曲面 $E(\theta, \phi; J)$ 在角度球面流形上展示了自旋 $J$ 约束下的微观能量分布。

当自旋低于临界值 $J_{\text{crit}}$ 时，势能极小位于主平面内（$\phi = 0^\circ$ 或 $90^\circ$），体系处于经典的共面进动相；当自旋超过 $J_{\text{crit}}$ 时，主平面鞍点发生超叉分岔（Pitchfork Bifurcation），能量极小漂移至非平面空间角度（$\theta_0 \approx 55^\circ, \phi_0 \approx 30^\circ$），在势能面上形成一对完全简并的左手与右手深阱。3D-TAC 是在微观密度泛函层次上揭示手征相变的理论核心。

\subsection{三轴投影壳模型（TPSM）}

为了兼具大规模壳模型的组态相互作用优势与形变基底的紧凑性，K. Hara、J. A. Sheikh 与孙扬等人发展了\term{三轴投影壳模型}（Triaxial Projected Shell Model, TPSM）。

TPSM 的核心思想是在具有三轴形变（$\beta_2, \gamma$）的变形 BCS 多准粒子真空基底 $|\Phi(\beta, \gamma)\rangle$ 上，利用三维不可约角动量投影算符（3D Angular Momentum Projection Operator）恢复空间转动对称性：
\begin{equation}
    \hat{P}_{MK}^I = \frac{2I+1}{8\pi^2} \int d\Omega \, D_{MK}^{I*}(\Omega) \hat{R}(\Omega)
    \label{eq:projection_operator}
\end{equation}
将投影出的纯净角动量基矢组 $|I M K, \kappa\rangle = \hat{P}_{MK}^I |\Phi_\kappa\rangle$ 进行组态混合对角化：
\begin{equation}
    \sum_{K' \kappa'} \left( H_{K \kappa, K' \kappa'}^I - E_I N_{K \kappa, K' \kappa'}^I \right) C_{K' \kappa'}^I = 0
    \label{eq:tpsm_secular}
\end{equation}
TPSM 不仅能够完美重现偶偶核与奇核低激发转动态与 $\gamma$ 振动带，更能高精度复现高自旋回弯与手征伙伴带的奇偶交错跃迁几率，且严格保持了实验室坐标系的角动量量子数守恒。

\subsection{模型无关四极形状不变量与电荷闭合定理}

在实验谱学中，如何从不依赖任何理论模型拟合假定的纯实验观测量中，提取原子核真实的内禀形状？Kumar、Cline 以及 Kumahora-Sugawara 引入了四极张量标量积不变量方法：
\begin{equation}
    q_2(I) = \frac{1}{2I+1} \sum_{I'} |\langle I || Q_2 || I' \rangle|^2 \propto \langle \beta^2 \rangle
\end{equation}
\begin{equation}
    q_3(I) = -\sqrt{\frac{7}{10}} \frac{\sqrt{35}}{10} \frac{1}{2I+1} \sum_{I' I''} \langle I || Q_2 || I' \rangle \langle I' || Q_2 || I'' \rangle \langle I'' || Q_2 || I \rangle W(2 2 I I''; 2 I') \propto \langle \beta^3 \cos 3\gamma \rangle
\end{equation}

由此得到绝对模型无关的有效三轴角：
\begin{equation}
    \cos 3\gamma_{\text{eff}} = \frac{q_3(I)}{[q_2(I)]^{3/2}}
\end{equation}

在近期对于奇钡（$^{129\text{--}137}\text{Ba}$）等核的电荷闭合核验中发现，如果只在单一自旋通道（例如仅针对基带 $I'=I+2$）进行求和截断，受限通道虽然汇聚了绝大部分电磁衰变率，但仅能代表全部四极电荷分布不变量的 $20\%\sim 30\%$；在截断空间中计算的有效角甚至会出现 $\cos 3\gamma > 1$ 的非物理虚假越域。必须完成全部多极自旋通道求和，电荷几何分布方能完全闭合。这一原则为全书后续数据解释奠定了不可动摇的方法论基石。
